\documentclass[a4paper,11pt]{article}

\usepackage[utf8]{inputenc}
\usepackage[T1]{fontenc}
\usepackage[polish]{babel}
\usepackage{amsmath, amssymb, amsthm}
\usepackage{listings}
\usepackage{xcolor}
\usepackage{geometry}

\geometry{margin=2cm}

\lstset{
    language=Python,
    basicstyle=\ttfamily\small,
    keywordstyle=\color{blue},
    stringstyle=\color{red},
    commentstyle=\color{green!50!black},
    numbers=left,
    numberstyle=\tiny\color{gray},
    frame=single,
    breaklines=true,
    captionpos=b
}

\title{Wybrane zagadnienia Algebry 2026\\Sprawozdanie 2.}
\author{Sara Żyndul \\ 279686}
\date{\today}

\begin{document}

\maketitle

\section{Laboratorium 4: Dzielenie wielomianów wielu zmiennych}

\subsection{Ad a) Struktura obsługująca wielomiany wielu zmiennych}
Wielomiany wielu zmiennych zostały zaimplementowane w oparciu o strukturę słownikową (tzw. reprezentację rzadką). Każdy wielomian to obiekt słownika (\textit{dictionary}), w którym:
\begin{itemize}
    \item \textbf{Klucz} reprezentuje wektor potęg poszczególnych zmiennych w postaci niemutowalnej krotki (\textit{tuple}). Na przykład dla pierścienia $\mathbb{R}[x,y,z]$ klucz \texttt{(2, 7, 0)} odpowiada jednomianowi $x^2y^7$.
    \item \textbf{Wartość} odpowiada niezerowemu współczynnikowi rzeczywistemu przy danym jednomianie.
\end{itemize}
Taka struktura pozwala uniknąć wybuchu pamięciowego (charakterystycznego dla reprezentacji gęstej) i ułatwia wykonywanie operacji algebraicznych – dodawanie wielomianów sprowadza się do sumowania wartości przypisanych do tych samych kluczy.

\subsection{Ad b) Porządki na jednomianach}
Wielomiany wielu zmiennych nie posiadają jednego, naturalnego wyrazu wiodącego (\textit{Leading Term} -- LT). Hierarchię tę należy narzucić poprzez określenie porządku:
\begin{itemize}
    \item \textbf{Standardowy porządek leksykograficzny (Lex):} Porównuje wektory potęg od lewej do prawej, co jest tożsame ze słownikowym uporządkowaniem wyrazów. Pierwsza zmienna (np. $x$) ma absolutny priorytet nad kolejnymi.
    \item \textbf{Permutowany Lex:} Działa identycznie jak standardowy porządek leksykograficzny, lecz przed porównaniem zmieniona zostaje kolejność zmiennych w wektorze. Pozwala to wymusić nienaturalną hierarchię, na przykład wymusić, aby zmienna $y$ dominowała nad $z$, a $z$ nad $x$ ($y > z > x$).
    \item \textbf{GradedLex (GrLex):} Porządek z gradacją (stopniowany). W pierwszej kolejności porównywany jest całkowity stopień jednomianu (suma wszystkich potęg zmiennych). Jednomian o wyższym stopniu całkowitym jest uznawany za większy. Dopiero w przypadku zrównania stopni, rozstrzygnięcie następuje na podstawie standardowego porządku Lex.
\end{itemize}

\subsection{Ad c) Funkcja \texttt{PolynomialReduce}}
Zaimplementowana funkcja \texttt{PolynomialReduce} realizuje algorytm dzielenia wielomianów wielu zmiennych przez ciąg dzielników. Poniżej przedstawiono kluczowy fragment kodu odpowiedzialny za pętlę redukującą:

\begin{lstlisting}[language=Python, caption={Kluczowy fragment pętli redukującej w PolynomialReduce}]
while p:
    lt_p_exp, lt_p_coeff = get_lt(p, order_key)
    division_occurred = False

    for i, g in enumerate(G):
        lt_g_exp, lt_g_coeff = get_lt(g, order_key)
        if not lt_g_exp: continue
        
        # Sprawdzenie czy LT(g_i) dzieli LT(p)
        if divides(lt_g_exp, lt_p_exp):
            mult_exp = tuple(e1 - e2 for e1, e2 in zip(lt_p_exp, lt_g_exp))
            mult_coeff = lt_p_coeff / lt_g_coeff
            
            # Aktualizacja i-tego ilorazu (alfa)
            alphas[i][mult_exp] = alphas[i].get(mult_exp, 0) + mult_coeff
            alphas[i] = clean_poly(alphas[i])
            
            # Odjecie od p wielokrotnosci dzielnika g_i
            term_times_g = term_mult(mult_coeff, mult_exp, g)
            neg_term = {k: -v for k, v in term_times_g.items()}
            p = poly_add(p, neg_term)
            
            division_occurred = True
            break  # Algorytm jest zachlanny
            
    if not division_occurred:
        # Nieredukowalny wiodacy jednomian trafia do ostatecznej reszty
        r[lt_p_exp] = r.get(lt_p_exp, 0) + lt_p_coeff
        r = clean_poly(r)
        del p[lt_p_exp]
\end{lstlisting}

\subsection{Ad d) Rozwiązanie Ćwiczenia 37 i wpływ kolejności}
Przeanalizowano porządek GradedLex ($x>y>z$) dla wielomianów: $f = x^3 - x^2y - x^2z$, $g_1 = x^2y - z$ oraz $g_2 = xy - 1$.
Zgodnie z implementacją uzyskano następujące wyniki redukcji:
\begin{itemize}
    \item Redukcja $f$ przez ciąg $(g_1, g_2)$ dała resztę: $r_1 = x^3 - x^2z - z$.
    \item Redukcja $f$ przez ciąg $(g_2, g_1)$ dała resztę: $r_2 = x^3 - x^2z - x$.
\end{itemize}

\textbf{Uzasadnienie różnicy:} 
Wyniki są różne ze względu na \textit{zachłanność} algorytmu dzielenia. W procesie redukcji dla wyjściowego wielomianu pojawia się składnik $-x^2y$. Okazuje się, że wiodące jednomiany z obu dzielników (zarówno $LT(g_1) = x^2y$, jak i $LT(g_2) = xy$) są w stanie bez reszty podzielić ten składnik. Algorytm nie dokonuje analizy głębszej i wybiera pierwszy pasujący dzielnik z podanego ciągu. Zamiana kolejności na $(g_2, g_1)$ sprawia, że składnik $-x^2y$ zostaje zredukowany przez $xy$ zamiast przez $x^2y$, co odcina od $f$ zupełnie inny "ogon" algebraiczny (odpowiednio $-z$ lub $-x$).

Takie zachowanie dowodzi, że ciąg wielomianów $(g_1, g_2)$ nie stanowi Bazy Gröbnera generowanego przez siebie ideału. Warto jednak odnotować fundamentalną własność algorytmu: różnica uzyskanych reszt $r_1 - r_2$ jest zawsze elementem ideału $\langle g_1, g_2 \rangle$.

\subsection{Ad e) Redukcja wielomianu $h(x,y,z)$ dla zmiennych z indeksu}
Na podstawie przydzielonych zmiennych $a=2, b=7, c=9, d=6, e=8, f=6$ zdefiniowano wielomian:
\begin{equation}
    h(x,y,z) = x^2 y^7 - y^9 z^6 + x^8 z^6
\end{equation}
Zaproponowano ciąg dzielników $G = (g_1, g_2)$, gdzie $g_1 = x^2 - z$ oraz $g_2 = y^7 - x$. 
Dokonano ewaluacji wywołania \texttt{PolynomialReduce(h, G)} dla trzech różnych porządków leksykograficznych (zadanych przez zmianę permutacji priorytetów zmiennych). Eksperyment zwrócił trzy zupełnie różne reszty:

\begin{lstlisting}[basicstyle=\ttfamily\small, frame=single, title={Wynik z terminala}]
Testowanie wielomianu h dla G = [x^2 - z, y^7 - x]:

Nowa reszta dla porzadku x > y > z:
Reszta (r): {(0, 9, 6): -1.0, (0, 7, 1): 1.0, (0, 0, 10): 1.0}

Nowa reszta dla porzadku y > x > z:
Reszta (r): {(1, 2, 6): -1.0, (1, 0, 1): 1.0, (0, 0, 10): 1.0}

Nowa reszta dla porzadku y > z > x:
Reszta (r): {(13, 2, 0): -1.0, (20, 0, 0): 1.0, (3, 0, 0): 1.0}
\end{lstlisting}

\textbf{Źródło obserwowanych różnic:}
Wybór określonej permutacji zmiennych fundamentalnie zmienia to, który ze składników staje się jednomianem wiodącym $LT$. Na przykład:
\begin{itemize}
    \item W porządku Lex przy zmiennych ($x > y > z$) wyrazem wiodącym $h$ jest $x^8 z^6$.
    \item W porządku Lex przy zmiennych ($y > x > z$) wyrazem wiodącym $h$ staje się $-y^9 z^6$.
\end{itemize}
Zmiana definicji $LT(h)$ (a także $LT(g_1)$ i $LT(g_2)$) całkowicie modyfikuje krok początkowy algorytmu. Redukcja wielomianu $h$ zaczyna się od "zgryzania" zupełnie innego składnika, co determinuje wygenerowanie innego zestawu składników podlegających redukcji w kolejnych iteracjach. W efekcie (ponieważ $G$ ponownie nie jest Bazą Gröbnera) ostateczne nieredukowalne ułamki trafiające do reszty przyjmują w każdej permutacji odmienną algebraiczną formę.
\end{document}